\vspace{-4mm}
\section{Introduction}
\vspace{-2mm}
The attention mechanism \cite{bahdanau2014neural,vaswani_attention_2017} has been shown to be an important primitive for accurate sequence modeling. Attention is moreover efficient during training as it is rich in matrix multiplications and  can thus take advantage of highly parallel processing capabilities and specialized accelerators on modern GPUs. However, the complexity of attention is quadratic in sequence length, and hence it is a fundamentally expensive primitive. And while recent techniques  have made it possible to scale attention to longer sequences through hardware-aware restructuring of the intermediate computations \cite{dao_flashattention_2022,dao_flashattention-2_2023,liu_ring_2023,brandon_striped_2023}, these methods still  require storing the key/value vectors of previous elements, and this  ``KV cache'' (whose size grows linearly) can  be unwieldy to manage for long sequences. 


Linear attention transformers \cite{katharopoulos_transformers_2020} replace the exponential kernel in softmax attention with a dot-product over (possibly transformed) key and query vectors. This   makes it possible to formulate linear attention as a linear RNN with matrix-valued hidden states, thus obviating the need for a KV cache and enabling constant-memory inference. While initial variants of linear attention generally underperformed softmax attention    on language modeling,  gated variants of linear attention which incorporate a data-dependent gating factor have  recently been shown to be competitive against strong transformer baselines \cite{yang_gated_2023,qin_hgrn2_2024,beck2024xlstm, peng_eagle_2024}. 
These gated linear transformers, along with time-varying state space models such as Mamba \cite{gu_mamba_2023,mamba2} (which can be reparameterized as a gated linear transformer \cite{yang_gated_2023}), have been suggested as a  potential alternative to ordinary transformers.
However, despite the competitive language modeling performance, these models have been shown to underperform transformers on recall-intensive tasks \cite{arora_simple_2024, arora-2024-jrt}, which is important for many practical downstream tasks of interest (e.g., in retrieval-augmented generation \cite{lewis_retrieval-augmented_2021}). 



To enhance  associative recall over long contexts, \citet{schlag_linear_2021} propose DeltaNet, a variant of a linear transformer which uses a  delta rule-like update   \cite{widrow_adaptive_1988} to retrieve  and update a value vector that is associated with the current key. DeltaNet was found to be effective on synthetic tasks and small scale language modeling/machine translation. However, the original work used a  sequential algorithm that did not parallelize across sequence length, thus resulting in hardware-inefficient training, and it has not been clear how to scale DeltaNet    to larger models and datasets.

This work describes a hardware-efficient training algorithm for DeltaNets which parallelizes the forward/backward passes across  sequence  length. 
We reparameterize the DeltaNet as a matrix-valued RNN whose recurrence is given by a generalized Householder transformation. This reparameterization enables the use of the compact WY representation \cite{bischof_wy_1985} for  products of Householder matrices, 
eliminating the need to materialize the hidden states of matrix size at each time step during parallel training, which would otherwise result in high I/O costs.  The memory-efficient representation makes it possible to straightforwardly extend the chunkwise parallel strategy for training linear attention models \cite{hua_transformer_2022,sun2023retentive,yang_gated_2023} to the DeltaNet case. We scale  DeltaNets to moderate-scale language modeling benchmarks (1.3B models trained on 100B tokens), where DeltaNet is found to obtain better language modeling and zero-shot downstream task performance than strong linear recurrent models such as Mamba \cite{gu_mamba_2023} and GLA \cite{yang_gated_2023}. For in-context retrieval and learning evaluation, we evaluate DeltaNet on synthetic and real benchmarks \cite{zoology,akyurek2024context,poli_mechanistic_2024, arora_simple_2024}, where it is again found to perform well against linear recurrent baselines. Finally, we experiment with a hybrid approach where we combine DeltaNet layers with sliding attention layers or global attention layers, and find that these hybrid models can improve upon ordinary transformers, as well as the pure DeltaNet transformer.

